
%%%%%%%%%%%%%%%%%%%%%%%%%%%%%%%%%%%%%%%%%%%%%%%%%%%%%%%%%%%%%
\subsubsection{具体分析}

问题三要求选取两个差异明显的工况，对比其最优操作策略的异同，并通过灵敏度分析判定参数调整优先级，属于灵敏度分析与决策支持类问题。

\textbf{两工况选取依据。}浓度是影响除尘负荷最直接的因素，选取最高浓度工况（工况0，$C_{in}=44.34\;\text{g/Nm}^3$）与最低浓度工况（工况5，$C_{in}=26.01\;\text{g/Nm}^3$）进行对比，最能体现策略差异。

\textbf{灵敏度分析方法选择。}$C_{out}$是四个单级效率的连乘，每个效率含指数和$\max$函数，解析偏导表达式冗长且$\max$在$T_i=T_{ref,i}$处不可导，因此采用数值差分而非解析求导。中心差分误差$O(h^2)$，前向差分误差$O(h)$，同等步长精度更高，故选用中心差分。

\textbf{优先级判定的量纲问题。}直接用$|S^C/S^P|$（浓度灵敏度与电耗灵敏度之比）有量纲问题：电压的比值单位是$(\text{mg/Nm}^3)/\text{kW}$，振打的比值单位不同，直接平均不合理。需改用无量纲弹性系数消除量纲差异，使电压与振打可比。
